\section{Discussion}
\label{sec:discussion}

The introduction asked four questions. The answers are partial, and the first and the third carry the most weight.

\emph{Pass interactions.} The instruction-count interactions are exact, large, and structured. Of the \InstInteractionCells{}
ordered pairs of the \NumThreePasses{} R3 passes, \InstInteractionConsistentCells{} have a mean beyond two standard errors across
kernels. The largest entries are self-pairs and pairs of passes that remove the same redundant code. Where a second application of a
pass leaves the bitcode unchanged, the self-interaction equals minus the pass's single effect, and that identity holds exactly in
all \InstSelfIdentityCases{} cases tested. The timing interactions add little to this picture. \InteractionSignificantCells{} of
\InteractionCells{} normalized timing interactions pass the two-standard-error test. That is more than chance alone would give, but
every entry is smaller than the noise floor, so we do not count them as established.

\emph{Canonical form.} On the \texttt{mem2reg}-only form, the five loop passes change every kernel. On the sweep canonical form, the
smallest count among them is \PassChecksLoopMinCanonical{}, for \texttt{loop-unroll}. The canonical form therefore matters for
\texttt{loop-unroll}, so the choice of canonical form is part of the result and not only a setup detail. For the other passes, the
two forms differ by at most \PassChecksOtherMaxDiff{} kernels.

\emph{Redundant executions.} Redundancy is common. On average, \RpbOverallMean{} of the pass executions in these sweeps leave the IR
unchanged. Charged with their share of the pass time, the redundant executions consume \WasteTotalSeconds{} seconds across all
sweeps. Only a small part of the optimizer's wall time is pass time, so the wall time is not used for the waste.

\emph{Pareto structure.} The fronts contain \ParetoPoints{} non-dominated points. Their hypervolume depends on the reference point,
so a value is meaningful only with its reference. Two of the three objectives depend on timing measurements, which carry noise. The speed objective uses
the timing, and the compile objective uses the optimizer wall time.

\emph{Verification.} The block-local checker decides refinement for the blocks in its fragment, which are \BlocksEncodablePct{}\% of
the matched block pairs. Within those blocks, \NamesRefine{} names refine and \NamesViolated{} are violated. In
\WitnessGainSweepsPct{}\% of the SAT sweeps, the candidate counterexample names an instruction to which the optimizer added a flag.
The remaining candidates involve other changes or upstream values. The benchmark outputs match the canonical outputs in all but
\OutputsDiffering{} of \OutputsChecked{} sweeps, on the one input each kernel runs. None of the exceptions is a SAT sweep. For each of them, the difference disappears under \texttt{-O2} lowering. That is consistent with the candidates being unreachable on those
inputs, but it is not a proof of correctness. The one candidate inspected in full is excluded by an invariant of its loop, not by any
single block, so SAT answers are reported as candidates for inspection, not as miscompilations.

Four implications follow. First, a measure that is exact for each kernel, such as the instruction count, should be preferred to a
noisy one whenever the question allows it. Second, the canonical form is a parameter of the study, and it should be reported and
tested, as Section~\ref{sec:results-canonical} does for the loop passes. Third, block-local verification of optimized code needs
facts from the surrounding function, such as value ranges, before its SAT answers can be trusted without inspection. Fourth,
benchmark output can depend on timing. A canonical control failed on the bitcount kernel, and the kernel's ranking lines were
masked under a rule recorded with the kernel. Any campaign that compares outputs should include such a control, and it should also
test the lowering level, as the exceptions of Section~\ref{sec:results} show.

Future work follows from these limits. The first item is a sensitivity analysis of the canonical form for every pass, not only the
loop passes. The second is a measurement method with less noise, so that the timing interactions can be measured below the noise
floor. The third is a function-level encoding with range facts, which would reduce the SAT answers that depend on context. The
fourth is a compile-time objective based on pass time, which would change the fronts of Section~\ref{sec:results}. Larger kernel
sets, a second machine, and whole-program benchmarks would test whether these results hold beyond small loop kernels.

\section{Conclusion}

We measured the interactions among the passes of the LLVM new pass manager on \NumKernels{} kernels, over \NumSweeps{} pipeline
variants. The main measure is the instruction count of the optimized IR, which is exact and does not pass through the backend. Of
its pairwise interactions, \InstInteractionConsistentCells{} of \InstInteractionCells{} are consistent across kernels, and the
self-interactions follow an identity that holds in every tested case. Timing resolves few of these interactions,
and the ones it flags lie below the noise floor. Redundant executions are common. The refinement checker decides refinement for a
minority of the blocks it examines, and in \WitnessGainSweepsPct{}\% of its SAT sweeps the optimizer added a flag to the instruction
that the SAT answer names. The one such case inspected in full is a false alarm of the block-local check. \OutputsDiffering{} sweeps
produce output that differs from the canonical output under the \texttt{-O0} lowering, and we have not explained the difference.
These are the claims the data support. They are bounded by one compiler release, one lowering level for the timing
results, and one input per kernel.

\subsection{Threats to validity}
\label{sec:threats}

\textbf{Measures.} The main results use the instruction count of the optimized IR. That count measures code size at the IR
level, and it says nothing about running time. The timing results carry measurement noise, and the speed objective of the
Pareto analysis inherits it. The compile objective uses the optimizer wall time, which includes process start-up and trace output.
Changing that objective to pass time would change the fronts, and the paper does not redo them.

\textbf{Lowering.} The timing results and the text sizes use an optimized module lowered with \texttt{llc -O0}, so they exclude backend optimization. The instruction counts do not depend on lowering. A pilot (Section~\ref{sec:methods}) compares the size change of four passes on four kernels at \texttt{-O2} and at \texttt{-O0}. It gives \LoweringPilotNonzeroOZero{} nonzero size changes at \texttt{-O0} out of \LoweringPilotCases{} cases, against \LoweringPilotNonzeroOTwo{} at \texttt{-O2}, and \LoweringPilotSignFlips{} cases change sign between the two levels. The \texttt{-O0} lowering was chosen to keep backend optimization out of the size and timing results. It has a cost of its own, because it can change program output. The four output exceptions of Section~\ref{sec:results} are sweeps of \texttt{susan-smoothing} whose \texttt{-O0} binaries differ from the canonical binary, and for each of them the same optimized IR lowered at \texttt{-O2} gives the canonical output. The pilot used object files, not linked executables, so its sizes are a proxy for the sweeps' binaries.

\textbf{Canonical form.} The sweeps use the sweep canonical form, which applies \texttt{mem2reg}, \texttt{loop-simplify}, and \texttt{lcssa}. Section~\ref{sec:results-canonical} shows that the loop passes depend on this choice. On the \texttt{mem2reg}-only form, every loop pass changes every kernel, whereas on the sweep canonical form \texttt{loop-unroll} changes only \PassChecksUnrollCanonicalChanged{} of \PassChecksKernels{} kernels. The other passes are less sensitive, and the two forms differ by at most \PassChecksOtherMaxDiff{} kernels for each of them. A pass that depends on a loop form it does not find may do nothing on these inputs, and the study does not separate that case from a pass that has no effect.

\textbf{Timing.} Paired windows remove the bias of unpaired baselines, and the identical-IR runs bound the noise of identical code. They do not bound the differences between distinct binaries. Code layout and alignment can change the running time of two binaries that compute the same thing, and the identical-IR runs cannot show that effect. For this reason the timing comparisons count exceedances of the noise floor but do not treat them as interference. \textbf{Benchmarks and inputs.} The \NumKernels{} kernels are small, and most of them are loop nests or control-flow kernels. Each kernel runs on one input. The output agreement reported in Section~\ref{sec:results} is therefore evidence about those inputs only, and a miscompilation that appears only on other inputs would not be detected. The bitcount kernel shows that benchmark output can depend on timing. Its ranking lines were masked after the control check failed, and the rule is recorded with the kernel. \OutputsDiffering{} sweeps of \texttt{susan-smoothing} produce output images that differ from the canonical images under the \texttt{-O0} lowering, so the output comparison is not a complete check, and the paper reports those sweeps as exceptions.

\textbf{One compiler release.} All results use LLVM \LLVMVersion. \RejectedCanonicalCount{} canonical passes are rejected by
this release, and \SubstituteCount{} substitutes replace them. An earlier campaign recorded \ArchivedCrashes{} crashes of one single-pass sweep type, \texttt{unify-loop-exits}, on specific kernels. The current campaign replaces that pass by \texttt{newgvn} in its slot on every kernel before any sweep runs. The size levels use forced function attributes, because the release
rejects the \texttt{Os} and \texttt{Oz} pipeline names. A standalone \texttt{instcombine} aborts on some functions. The spelling that skips the fixpoint check gave identical IR on \SpellingIdentical{} of the \SpellingCandidates{} completed sweeps that contain the pass, and \SpellingDiffering{} sweeps differed. It is still a different spelling of the same pass. All of these choices could change the results in another release.

\textbf{Scope of the refinement proofs.} The UNSAT answers are proofs for block pairs in the fragment of
Section~\ref{sec:unsat}, for names defined in both blocks, and for non-poison live-ins. They say nothing about memory, calls,
floating point, vectors, control flow inside the block, or facts that hold only in the surrounding function. Only
\BlocksEncodablePct{}\% of matched block pairs lie in the fragment, so most of each program is outside the proofs. A SAT answer is a
candidate counterexample, and the paper reports it as one. The encoder's agreement with the reference semantics is checked
exhaustively for small widths, which is finite evidence for the lemma of Section~\ref{sec:framework}.

\textbf{Separate evidence.} The timing measurements, the refinement checks, and the output comparisons are separate evidence. A
sweep whose optimized IR refines its canonical IR can still be slower, and a faster sweep is not thereby correct. The paper draws
conclusions about correctness only from the refinement results and the output comparisons, and about cost only from the
instruction counts and the timings.

\textbf{Estimators.} The single-effect ratio baseline is a diagnostic only. Its denominator is noise-sized for
\MpassNoiseDominatedPct{}\% of the kernel-entry pairs, and the regime R6 chains are selected from the instruction-count matrix instead, so they do not inherit that instability. The pass-time waste assumes that every execution costs the same share of its sweep's pass time. The hypervolume depends
on the reference point, as Figure~\ref{fig:hv} shows, and the canonical reference was fixed before the analysis.


